\chapter{Cel i zakres pracy}
\label{chap:cel}

% Rozdział powinien zmieścić się na jednej stronie.

Rozdział otwiera się jednoznacznym sformułowaniem celu pracy, czyli tego, co
w~jej ramach zostało zaprojektowane, obliczone, zbadane lub ocenione.
Przykładowo: \emph{Celem pracy jest opracowanie koncepcji instalacji
grzewczej budynku jednorodzinnego zasilanej z~powietrznej pompy ciepła oraz
porównanie tego rozwiązania z~kotłem gazowym pod względem kosztów
eksploatacji i~emisji CO\textsubscript{2}.}

Następnie wymienia się zadania, których realizacja doprowadziła do
osiągnięcia celu. Zakres pracy obejmował:
\begin{itemize}
  \item przegląd literatury oraz przepisów i~norm dotyczących
        projektowanego systemu,
  \item obliczenie projektowego obciążenia cieplnego budynku zgodnie
        z~normą PN-EN 12831-1 \cite{pn12831},
  \item dobór urządzeń oraz obliczenia hydrauliczne instalacji,
  \item analizę ekonomiczną i~środowiskową rozpatrywanych wariantów,
  \item sformułowanie wniosków i~zaleceń projektowych.
\end{itemize}

W~tej części warto zapowiedzieć układ pracy, np.: \emph{W~rozdziale drugim
omówiono stan wiedzy na temat \dots, rozdział trzeci zawiera obliczenia
\dots, a~w~rozdziale czwartym przedstawiono \dots} Należy tu również
wymienić programy komputerowe wykorzystane do obliczeń, symulacji
i~przygotowania rysunków, np.\ QGIS \cite{qgis} albo EPANET \cite{epanet22}.
